\section{Conducción autónoma, robótica y cómo expresar una tesis en el mercado privado y en el mercado público}

Las secciones anteriores trataron las empresas de modelos y la tecnología financiera; esta sección agrega la dirección del mundo físico: la conducción autónoma y la robótica. Freda valora mucho la capacidad de ejecución de Waymo: lanza el servicio en nuevas ciudades cada vez más rápido, amplía sus zonas de operación y sus ingresos empiezan a poder estimarse. Al mismo tiempo, advierte que las empresas de robótica alcanzan con facilidad valuaciones muy altas en el mercado privado (primario), mientras que en el mercado público (secundario) puede ser más sencillo expresar la tendencia mediante Nvidia, Tesla, la cadena de suministro relacionada u otros valores cotizados.

\subsection{Waymo y la comercialización de la conducción autónoma}

En el episodio se menciona que Waymo opera en varias zonas y que el tiempo para lanzar el servicio en una ciudad pasó de varios años en San Francisco a plazos más cortos en Austin y Silicon Valley; además, el costo por ciudad no es tan alto como se pensaba. Lo relevante de este dato es que, una vez que la conducción autónoma pasa de la «demostración tecnológica» a la «expansión operativa», los inversionistas empiezan a calcular el costo de replicar el servicio en cada ciudad, el número de vehículos, los ingresos anualizados, la velocidad de aprobación regulatoria y la economía unitaria.

\begin{knowledgebox}{Variables que hay que observar en la comercialización de la conducción autónoma}
\begin{tabularx}{\textwidth}{p{0.23\textwidth}p{0.33\textwidth}X}
\toprule
Variable & Pregunta que hay que responder & Por qué importa \\
\midrule
Velocidad de apertura de ciudades & Cuánto tarda una nueva ciudad en pasar de pruebas a operación & Determina la velocidad con que se replican los ingresos. \\
Costo por ciudad & Cuánto capital e inversión operativa requiere entrar en una ciudad & Determina si la expansión es sostenible. \\
Tamaño de la flota & Cuántos vehículos prestan servicio estable en las calles & Determina la oferta y la experiencia del usuario. \\
Permisos regulatorios & Si puede circular por autopistas, operar entre zonas y ampliar el área de servicio & Determina los límites del negocio. \\
Economía unitaria & Cuáles son los ingresos, el mantenimiento, el seguro y la depreciación por vehículo y por día & Determina el modelo de rentabilidad a largo plazo. \\
\bottomrule
\end{tabularx}
\end{knowledgebox}

\subsection{Robótica y su expresión en el mercado público}

Esta sección vuelve a situar la conducción autónoma y la robótica en la perspectiva de un fondo Crossover. Antes se dijo que la incertidumbre en el mercado privado es alta; por eso los inversionistas suelen preguntarse: ¿existe una forma más líquida y más transversal de expresar la misma tesis sobre la industria en el mercado público?

La robótica y la inteligencia corporeizada siguen siendo una dirección de largo plazo, pero en el mercado privado la incertidumbre de cada empresa individual es muy alta. Desde la perspectiva Crossover de Freda, la pregunta es: si crees en la robótica o en la infraestructura de IA, no necesariamente tienes que apostar por una sola empresa privada; también puedes expresar esa tesis mediante Nvidia, los chips, la energía, la nube, la cadena de suministro o plataformas que cotizan en bolsa. Este razonamiento se parece al de comprar Nvidia en 2023 para expresar la demanda agregada de las empresas de modelos.

\begin{warningbox}{El mercado privado y el mercado público expresan riesgos distintos}
En el mercado privado se puede obtener el rendimiento explosivo de una sola empresa, pero el capital queda inmovilizado por mucho tiempo, la valuación es poco transparente y la tasa de fracaso es alta. El mercado público tiene más liquidez y permite expresar la línea principal de una industria, pero también es más sensible a los factores macroeconómicos y al sentimiento del mercado. No se trata de que uno sea mejor que otro, sino de que tienen estructuras de riesgo y de horizonte temporal diferentes.
\end{warningbox}

\subsection{El mercado privado depende del consenso; el mercado público, de lo que no es consenso}

A continuación se resumen las diferencias entre el mercado privado y el mercado público. Los casos de Robinhood y Nvidia mostraron que el momento más interesante del mercado público suele llegar cuando el cambio real en la empresa ya ocurrió, pero el mercado todavía no lo ha incorporado por completo en el precio.

En el episodio hay una frase muy concisa: «el mercado privado depende del consenso; el mercado público, de lo que no es consenso». En el mercado privado, las empresas más solicitadas suelen necesitar consenso para levantar grandes rondas de financiamiento; en el mercado público, si todos están de acuerdo, el precio normalmente ya refleja ese consenso. El rendimiento excedente real proviene de entender el cambio en un negocio antes que el mercado, por ejemplo, la capacidad de ejecución de Robinhood y su expansión hacia varias líneas de negocio.

\subsection{Resumen de la sección}

La enseñanza común de la sección sobre conducción autónoma y robótica es que, en la tecnología de frontera, no basta con ver la demo: también hay que analizar las variables de comercialización y la forma de expresar la tesis. Un inversionista puede apostar por una empresa, o bien por la cadena de suministro, los chips, la nube, la energía o valores cotizados más líquidos del mercado público.
